\chapter{复分析：保角、保距与复积分选谱}
\label{chap:complex}

把实平面上的点 $(x,y)$ 记成复数 $z=x+iy$，两点距离就是 $|z-w|$。于是，旋转与反射都可以用复数公式计算。复分析还提供了另一种看谱定理的方式：把矩阵的特征值当作复平面上的点，用一条闭曲线圈出想要的谱成分。读本章时先把“是否保距”理解成检验 $|f(z)-f(w)|=|z-w|$，再看复可微性和复积分各自能说明什么。

\section{用复变函数辨别平面变换}

复导数 $f'(z)$ 由差商 $[f(z+h)-f(z)]/h$ 在复数 $h$ 趋于零时的极限给出。与实数情形不同，$h$ 可沿平面中的任意方向趋于零，极限必须相同。一个函数在某开集的每一点都复可微，就称它在该开集上全纯。

\begin{example}{用复数写出椭圆主轴的旋转}
设 $Q=\frac1{\sqrt5}\begin{pmatrix}2&-1\\1&2\end{pmatrix}$。把 $Q$ 写成复平面上的映射，求 $1+i$ 的像，并证明这个映射保距。
\end{example}
\begin{solution}
对 $z=x+iy$，直接相乘：
\[
 \frac{2+i}{\sqrt5}z=\frac{2x-y}{\sqrt5}+i\frac{x+2y}{\sqrt5}.
\]
实部、虚部恰好是 $Q(x,y)^T$ 的两个分量。系数 $a=(2+i)/\sqrt5$ 的模为 $1$，辐角为 $\arctan(1/2)$，故这是逆时针旋转。$1+i$ 的像为 $(1+3i)/\sqrt5$。任取两点 $z,w$，
\[
 |az-aw|=|a|\,|z-w|=|z-w|.
\]
因此一对点的距离和整条曲线的形状大小都保持。再加一个固定复数 $b$ 也保距，因为计算差时 $b$ 抵消。
\end{solution}

\begin{example}{关于直线的反射为何不是全纯函数}
把点关于直线 $y=x$ 作反射。写出复数形式，计算 $2+i$ 的像，并判断它在何处复可微。
\end{example}
\begin{solution}
反射交换两个坐标，因此 $R(x+iy)=y+ix=i\overline{x+iy}$，$R(2+i)=1+2i$。对任意两点，
\[
 |R(z)-R(w)|=|i\overline{z-w}|=|z-w|.
\]
但在任一点 $z$，差商为 $i\bar h/h$。取非零实数 $h$ 趋于 $0$，差商恒为 $i$；取纯虚数 $h=it$ 趋于 $0$，差商恒为 $-i$。两个方向的极限不同，所以处处不复可微。由于其共轭函数 $\overline{R(z)}=-iz$ 全纯，$R$ 称为反全纯映射。一般而言，关于过原点、与实轴夹角为 $\theta$ 的直线反射，可以先旋转 $-\theta$、再共轭、最后旋转 $\theta$，所得公式为 $e^{2i\theta}\bar z$。
\end{solution}

\begin{example}{局部保角为何不推出保距}
考察 $f(z)=z^2$ 在 $z_0=1+i$ 的作用。比较相互垂直的切向方向 $1$、$i$ 的像，并说明此映射能否保持任意两点距离。
\end{example}
\begin{solution}
$f(z_0+h)-f(z_0)=2z_0h+h^2$，故一阶近似为乘以 $2+2i$。两个切向量分别变成 $2+2i$ 与 $-2+2i$；对应实向量的内积为 $2(-2)+2\cdot2=0$。它们都逆时针旋转 $\pi/4$，长度乘 $2\sqrt2$。一般地，全纯函数在 $f'(z_0)\ne0$ 处，切向量都乘同一个非零复数 $f'(z_0)$，因而保持相交曲线的有向夹角。这是局部共形性。相比之下，$1,2$ 的距离为 $1$，它们的像 $1,4$ 的距离为 $3$，故 $f$ 不是欧氏等距映射。在 $z=0$，$f'(0)=0$，也不能套用上述保角论证。
\end{solution}

\begin{example}{把单位圆盘映回自身，也可能改变距离}
设 $a=\frac12$，$g(z)=\frac{z-a}{1-az}$。证明 $g$ 把开单位圆盘 $\{z:|z|<1\}$ 一一映到自身，且全纯、保角。检验它是否保持欧氏距离。
\end{example}
\begin{solution}
分母在圆盘内不为零。因 $a$ 是实数，两个模的平方分别为 $1-a(z+\bar z)+a^2|z|^2$ 与 $|z|^2-a(z+\bar z)+a^2$，相减得
\[
 |1-az|^2-|z-a|^2=(1-a^2)(1-|z|^2)>0.
\]
因此 $|g(z)|<1$。从 $w=g(z)$ 解得 $z=(w+a)/(1+aw)$；把 $a$ 换成 $-a$，同一个恒等式证明逆映射也把圆盘映入自身，所以 $g$ 是双射。它的导数
\[
 g'(z)=\frac{1-a^2}{(1-az)^2}
\]
处处非零，故在圆盘内保角。可是 $g(0)=-1/2$、$g(-1/2)=-4/5$，对应距离由 $1/2$ 变成 $3/10$。集合映到自身、局部保角与保持任意两点距离，是三件需要分别验证的事。
\end{solution}

\section{用柯西积分选出特征方向}

在以下两例中，$A=\begin{pmatrix}5&2\\2&2\end{pmatrix}$ 作用在 $\mathbb C^2$ 上。把实矩阵看成复矩阵之后，$Q^*=Q^T$（星号表示共轭转置），且 $A=Q\operatorname{diag}(6,1)Q^*$ 仍然成立。对于 $z\notin\{1,6\}$，矩阵 $(zI-A)^{-1}$ 称为预解式。这里 $z$ 是用于圈选特征值的复参数，向量本身则有两个复坐标。

沿逆时针简单闭曲线积分时，矩阵积分按每个元素分别计算。柯西积分公式告诉我们：若 $f$ 在曲线及其内部的一个开邻域上全纯，则 $f(z)/(z-\lambda)$ 的积分除以 $2\pi i$，在 $\lambda$ 位于曲线内部时等于 $f(\lambda)$，在 $\lambda$ 位于外部时等于 $0$。这里假定 $\lambda$ 不在曲线上。把预解式按特征值拆成分式，每个分式就有一个对应的极点，积分通过“圈内有无该极点”把这些分量区分开。

\begin{example}{绕特征值 $6$ 积分，得到秩一投影}
令 $\Gamma$ 为逆时针圆周 $|z-6|=2$。求
\[
 P=\frac{1}{2\pi i}\int_\Gamma(zI-A)^{-1}\,dz,
\]
并判定 $P$ 把向量投影到哪条直线上。
\end{example}
\begin{solution}
直接求逆得到
\[
 (zI-A)^{-1}=\frac{1}{(z-6)(z-1)}
 \begin{pmatrix}z-2&2\\2&z-5\end{pmatrix}.
\]
圆周圈住 $6$ 而不圈住 $1$，矩阵 $\frac1{z-1}\begin{pmatrix}z-2&2\\2&z-5\end{pmatrix}$ 的每个元素在圆盘上全纯。因此柯西积分公式直接把其中的 $z$ 换成 $6$：
\[
 P=\frac15\begin{pmatrix}4&2\\2&1\end{pmatrix}=qq^*,\qquad q=\frac1{\sqrt5}\binom{2}{1}.
\]
$q^*q=1$，故 $P^*=P$、$P^2=P$。$P$ 将 $x$ 变成 $q(q^*x)$，即沿着 $q$ 保留分量，垂直于 $q$ 的分量被抹去。又 $Aq=6q$，所以它是到特征值 $6$ 的特征直线的正交投影，记作 $P_6$。若换成圈住 $1$ 而不圈住 $6$ 的圆周，计算时将分子矩阵除以 $z-6$，再代入 $z=1$，得 $P_1=\frac15\begin{pmatrix}1&-2\\-2&4\end{pmatrix}=I-P_6$。积分所选择的是整个特征子空间，并不依赖于 $q$ 的相位选择。
\end{solution}

\begin{example}{复指数把自伴矩阵变成保距运动}
令 $t\in\mathbb R$，按指数级数定义 $e^{itA}$。用柯西积分写出它，证明它是幺正矩阵，并计算 $t=\pi/5$ 时的结果。
\end{example}
\begin{solution}
由于 $A=6P_6+P_1$，$P_6P_1=P_1P_6=0$，且两者都是幂等矩阵，所以 $A^n=6^nP_6+P_1$ 对所有 $n\ge1$ 成立；$A^0=I=P_6+P_1$。在收敛的指数级数中逐项代入，可得
\[
 e^{itA}=e^{6it}P_6+e^{it}P_1.
\]
同样地，把预解式写成 $P_6/(z-6)+P_1/(z-1)$，对正向圆周 $|z|=7$ 使用柯西公式，得到
\[
 e^{itA}=\frac{1}{2\pi i}\int_{|z|=7}e^{itz}(zI-A)^{-1}\,dz.
\]
由于 $|e^{6it}|=|e^{it}|=1$，两块分别只改变相位，于是 $(e^{itA})^*e^{itA}=P_6+P_1=I$。它保持 $\mathbb C^2$ 中任意两点的距离。取 $t=\pi/5$，$e^{6it}=-e^{it}$，故
\[
 e^{i\pi A/5}=e^{i\pi/5}(I-2P_6)
 =\frac{e^{i\pi/5}}5\begin{pmatrix}-3&-4\\-4&3\end{pmatrix}.
\]
右边的实矩阵因子 $I-2P_6$ 在特征方向 $q$ 上变号，在其垂直方向上保持不动；额外的相位因子也是等距的。原矩阵 $A$ 沿 $q$ 把长度扩大六倍，但经过函数 $\lambda\mapsto e^{it\lambda}$ 后，两块的伸缩模长都变成 $1$。
\end{solution}
